% Figure for Calculus §56.
\begin{tikzpicture}
  \begin{axis}[nfig, width=8.4cm, height=5.0cm, xmin=-0.8, xmax=11.5, ymin=0, ymax=0.19,
      xtick={0,2,4,6,8,10}, ytick={0.05,0.1,0.15},
      yticklabel style={/pgf/number format/fixed},
      xlabel={$x$}, ylabel={$y$}, xlabel style={at={(axis description cs:1,0)}, anchor=west},
      ylabel style={at={(axis description cs:0,1)}, anchor=south, rotate=-90}]
    \addplot[draw=none, fill=figB!18, domain=4:8, samples=60] {0.006*x*(10-x)} \closedcycle;
    \addplot[figB, domain=0:10, samples=100] {0.006*x*(10-x)};
    \addplot[figB, domain=-0.8:0] {0};
    \addplot[figB, domain=10:11.3] {0};
    \draw[figB!70, thin] (axis cs:4,0) -- (axis cs:4,0.144);
    \draw[figB!70, thin] (axis cs:8,0) -- (axis cs:8,0.096);
    \node[font=\small] at (axis cs:6,0.07) {$0.544$};
    \node[figB, font=\small, anchor=west] at (axis cs:7.9,0.145) {$y = f(x)$};
    \node[font=\scriptsize, align=center] at (axis cs:2.2,0.168) {total area under\\the graph $=1$};
  \end{axis}
\end{tikzpicture}
